\documentclass[12pt]{article}
\usepackage[margin=1in]{geometry}
\usepackage{amsmath,amssymb,amsthm}
\usepackage{mathtools}
\usepackage{enumitem}
\usepackage{hyperref}
\usepackage{xcolor}

\newtheorem{theorem}{Theorem}
\newtheorem{lemma}[theorem]{Lemma}
\newtheorem{proposition}[theorem]{Proposition}
\newtheorem{corollary}[theorem]{Corollary}
\newtheorem{definition}[theorem]{Definition}
\newtheorem{remark}[theorem]{Remark}
\newtheorem{assumption}{Assumption}

\newcommand{\R}{\mathbb{R}}
\newcommand{\E}{\mathbb{E}}
\newcommand{\supp}{\operatorname{supp}}
\newcommand{\argmax}{\operatorname{arg\,max}}
\newcommand{\argmin}{\operatorname{arg\,min}}

\title{Alternative Proof of Theorem~2 (Robust Trust):\\
Formalization via First-Order Conditions and the Envelope Theorem}
\author{Formalization report prepared for Piotr Dworczak\thanks{%
This document formalizes and verifies the alternative proof sketch
for Theorem~2 of Dworczak \& Smolin (2026), ``Robust Trust.''
The formalization was carried out using the MathPipeProver
soft-scaffolding pipeline (Claude Code orchestrating ChatGPT~5.4 Pro Extended).}}
\date{March 15, 2026}

\begin{document}
\maketitle

\begin{abstract}
We formalize the alternative proof of Theorem~2 (Robustly Rationalizable Solution) from Dworczak \& Smolin (2026). The original sketch proposes replacing the paper's Sion-based minimax argument with a first-order-conditions (FOC) plus envelope-theorem approach.
Our formalization confirms that this alternative route is \textbf{correct for the finite case} ($M$ and $\Theta$ finite), with three qualifications:
(i)~the envelope theorem must be applied via one-sided directional derivatives and a jointly chosen Bayes-optimal selector family, not via ordinary derivatives at an arbitrary optimizer;
(ii)~no new assumptions beyond the paper's finite-case setup are needed;
(iii)~the extension to infinite message spaces remains open.
The key technical insight beyond the sketch is a convex-analytic argument that extracts the selector family from a subgradient of the value function at the adversary's optimum.
\end{abstract}

\tableofcontents

%% =====================================================================
\section{Executive Summary}
%% =====================================================================

The sketch proposes a finite-case proof of Theorem~2 structured as follows:
\begin{enumerate}
\item Write down the agent's payoff under Bayes-optimal response to the adversary's strategy.
\item Derive the first-order conditions (FOCs) for the adversary's optimization problem.
\item Use an envelope theorem to eliminate the indirect effect (the dependence of the agent's optimal action on the perturbation of the adversary's strategy).
\item Show that the resulting FOCs imply the adversary's strategy remains optimal under agent commitment.
\item Conclude the minimax / saddle-point property.
\end{enumerate}

\paragraph{What was correct in the sketch.}
The posterior formula, the payoff decomposition, the derivative computation, and the commitment-game argument are all algebraically correct after notation alignment. (We use ``agent'' for the paper's decision maker throughout; the sketch calls this role ``DM.'')

\paragraph{What required repair.}
\begin{itemize}
\item \textbf{Envelope theorem precision.} The sketch invokes the envelope theorem informally. The correct tool is Milgrom--Segal (2002, \textit{Econometrica}), Theorem~3 (local equidifferentiability). This theorem gives \emph{one-sided} directional derivatives, not ordinary derivatives at an arbitrary optimizer. When multiple Bayes-optimal private strategies coexist, the sketch's claimed ordinary derivative $V'(0)$ is too strong.
\item \textbf{Selector family.} The fix (our key contribution beyond the sketch) uses a convex-analytic argument: at the adversary's optimum $\beta^*$, the subdifferential $\partial V(\beta^*)$ is nonempty by convex optimization on a compact polyhedron. A subgradient of $V$ at $\beta^*$ yields a \emph{jointly} chosen Bayes-optimal selector family $(\hat\sigma^*_m)_{m \in M}$ that satisfies the equal-payoff-on-support condition for all $\mu$ simultaneously. No uniqueness assumption is needed.
\item \textbf{Boundary cases.} The sketch differentiates as if the simplex constraint were absent. The correct treatment uses either a variational inequality on $\Delta(M)$ or KKT complementary slackness.
\end{itemize}

\paragraph{What remains open.} The extension to infinite $M$ (Block~G in the sketch) is explicitly left open by the author and remains unresolved.

%% =====================================================================
\section{Notation}
\label{sec:notation}
%% =====================================================================

We use the paper's notation throughout. The sketch uses different symbols for several objects; we record the correspondence here.

\medskip
\begin{center}
\begin{tabular}{lll}
\hline
\textbf{Sketch} & \textbf{Paper (canonical)} & \textbf{Object} \\
\hline
$\pi(m|\mu)$ & $\beta(m|\mu)$ & Misaligned adviser's reporting rule \\
$P(m)$ [sketch] & $q_\beta(m)$ & Marginal probability of message $m$ \\
$\gamma_m$ & $P_\beta(\cdot|m)$ & Agent's posterior given message $m$ \\
$a^*(\mu)$ & $\hat\sigma^*(\mu)$ & Bayes-optimal response at belief $\mu$ \\
$U(\pi, \sigma)$ & $U(\beta, \sigma)$ & Stage-game payoff \\
bad AI & misaligned adviser & Strategic adviser (prob.\ $1-\alpha$) \\
DM & agent & Decision maker / receiver \\
\hline
\end{tabular}
\end{center}

\begin{remark}[Notation collision]
The sketch writes $P(m)$ for the marginal message probability. In the paper, $P(m)$ already denotes the trust-region projection map (Definition~1 in the paper). We write $q_\beta(m)$ throughout to avoid collision.
\end{remark}

\begin{remark}[Suppressed objects]
The sketch suppresses the agent's private type $\theta$, the type-signal function $f$, the truthful strategy, and the paper's private-strategy notation $\hat\sigma$. We work in the simplified setting $u(a,\omega)$ with $|\Theta|=1$ for the main derivation, and note where $\theta$-dependence would require modification (see Remark~\ref{rem:theta}).
\end{remark}

%% =====================================================================
\section{Setup (Finite Case)}
\label{sec:setup}
%% =====================================================================

Fix a finite state space $\Omega = \{\omega_1, \ldots, \omega_N\}$ with full-support prior $\mu_0 \in \Delta(\Omega)$. Let $\tau \in \Delta(\Delta(\Omega))$ be the unconditional distribution of the adviser's posteriors, with finite support $M = \supp(\tau)$, $|M| = K < \infty$. Let $\alpha \in (0,1)$ be the alignment probability.\footnote{The restriction $\alpha > 0$ is used throughout to ensure $q_\beta(m) > 0$ for all $m \in M$. The case $\alpha = 0$ would require restricting to on-path messages and is excluded from this analysis.}

Let $A$ denote the agent's action set (compact metric) and $\Theta$ the finite set of private types, both as in the paper.

The misaligned adviser's strategy is $\beta : M \to \Delta(M)$, mapping each true posterior $\mu \in M$ to a distribution over reported messages. By the paper's WLOG argument (the adviser can always be taken to use only messages in $M$), we restrict to messages in $M$. Write $\mathcal{B} \coloneqq \prod_{\mu \in M} \Delta(M)$ for the adversary's strategy space and $\Sigma$ for the agent's strategy space (as defined in the paper).

\begin{proposition}[Posterior formula --- verified]
\label{prop:posterior}
For any $\beta$ and $m \in M$, the agent's posterior is
\begin{equation}
\label{eq:posterior}
P_\beta(\omega | m) = \frac{\alpha\, \tau(m)\, m(\omega) + (1-\alpha) \sum_{\mu \in M} \tau(\mu)\, \beta(m|\mu)\, \mu(\omega)}{q_\beta(m)},
\end{equation}
where
\begin{equation}
\label{eq:qbeta}
q_\beta(m) = \alpha\, \tau(m) + (1-\alpha) \sum_{\mu \in M} \tau(\mu)\, \beta(m|\mu).
\end{equation}
Since $\alpha > 0$ and $m \in M = \supp(\tau)$, we have $q_\beta(m) \geq \alpha\,\tau(m) > 0$ for all $m \in M$.
\end{proposition}

\begin{proof}
By Bayes' rule, $P_\beta(\omega|m) = \Pr_\beta(\omega,m) / \sum_{\omega'} \Pr_\beta(\omega',m)$. We decompose the joint probability $\Pr_\beta(\omega,m)$ into aligned and misaligned contributions.

\emph{Aligned contribution.} With probability $\alpha$, the adviser is aligned and reports truthfully. The joint probability that the state is $\omega$ and the aligned adviser sends message $m$ is $\alpha\,\tau(m)\,m(\omega)$, using the identity $\mu_0(\omega)\Pr(s=m|\omega) = \tau(m)\,m(\omega)$ (a standard consequence of the definition of posteriors).

\emph{Misaligned contribution.} With probability $1-\alpha$, the adviser is misaligned. Conditional on the adviser's true posterior being $\mu$, the probability of reporting $m$ is $\beta(m|\mu)$, and the probability of state $\omega$ given posterior $\mu$ is $\mu(\omega)$. Averaging over $\mu \sim \tau$ gives $(1-\alpha)\sum_{\mu \in M} \tau(\mu)\,\beta(m|\mu)\,\mu(\omega)$.

Adding the two contributions gives $\Pr_\beta(\omega,m) = \alpha\,\tau(m)\,m(\omega) + (1-\alpha)\sum_\mu \tau(\mu)\,\beta(m|\mu)\,\mu(\omega)$, which is the numerator in \eqref{eq:posterior}. Summing over $\omega$ gives $q_\beta(m)$.
\end{proof}

%% =====================================================================
\section{DM Payoff and First-Order Conditions}
\label{sec:foc}
%% =====================================================================

\subsection{The adversary's problem}

Throughout Sections~4--8, we write the derivation in the paper's private-strategy notation: $\bar{u}(\hat\sigma,\omega) \coloneqq \E[u(a,\omega,\theta) \mid \omega, \hat\sigma]$ denotes the state-contingent expected payoff from private strategy $\hat\sigma$, and $U(\hat\sigma,\nu) \coloneqq \sum_\omega \bar{u}(\hat\sigma,\omega)\,\nu(\omega)$ is the agent's payoff from $\hat\sigma$ at belief $\nu$. When $|\Theta|=1$, this reduces to the sketch's $u(a,\omega)$ notation.

For a fixed strategy $\beta$ of the misaligned adviser, the agent's expected payoff under Bayes-optimal response is
\begin{equation}
\label{eq:value}
V(\beta) = \sum_{m \in M} q_\beta(m)\, W\!\left(P_\beta(\cdot|m)\right),
\end{equation}
where $W(\nu) \coloneqq \max_{\hat\sigma} \sum_\omega \bar{u}(\hat\sigma, \omega)\,\nu(\omega)$ is the indirect utility at belief $\nu$, and $\bar{u}(\hat\sigma, \omega) \coloneqq \E[u(a,\omega,\theta) \mid \omega, \hat\sigma]$.

The adversary solves
\[
\min_{\beta \in \prod_{\mu \in M} \Delta(M)} V(\beta).
\]

\subsection{Perturbation and derivative}

Fix $\mu^* \in M$ and $m_1, m_2 \in M$. Define the simplex-preserving perturbation
\[
\beta_\varepsilon(m|\mu) = \begin{cases}
\beta(m|\mu^*) + \varepsilon\,[\delta_{m_1}(m) - \delta_{m_2}(m)] & \text{if } \mu = \mu^*, \\
\beta(m|\mu) & \text{otherwise},
\end{cases}
\]
which is feasible for $\varepsilon \in [-\beta(m_1|\mu^*),\, \beta(m_2|\mu^*)]$.

Define the unnormalized numerator $N_{\beta}(m,\omega) \coloneqq q_\beta(m)\,P_\beta(\omega|m)$. Under $\beta_\varepsilon$, only $N_{\beta_\varepsilon}(m_1,\omega)$ and $N_{\beta_\varepsilon}(m_2,\omega)$ change, and they change \emph{affinely} in $\varepsilon$:
\begin{align}
\frac{d}{d\varepsilon} N_{\beta_\varepsilon}(m_1,\omega)\Big|_{\varepsilon=0}
&= (1-\alpha)\,\tau(\mu^*)\,\mu^*(\omega), \label{eq:dN1} \\
\frac{d}{d\varepsilon} N_{\beta_\varepsilon}(m_2,\omega)\Big|_{\varepsilon=0}
&= -(1-\alpha)\,\tau(\mu^*)\,\mu^*(\omega). \label{eq:dN2}
\end{align}

Similarly, $q_{\beta_\varepsilon}(m_1)$ and $q_{\beta_\varepsilon}(m_2)$ are affine in $\varepsilon$ with slopes $\pm(1-\alpha)\tau(\mu^*)$.

\subsection{Direct and indirect effects}

Write $V(\beta) = \sum_m \Phi_m(\beta)$ where
\[
\Phi_m(\varepsilon) \coloneqq q_{\beta_\varepsilon}(m)\, W\!\left(P_{\beta_\varepsilon}(\cdot|m)\right)
= \max_{\hat\sigma} \sum_\omega \bar{u}(\hat\sigma,\omega)\, N_{\beta_\varepsilon}(m,\omega).
\]

For $m \notin \{m_1, m_2\}$, neither $N_{\beta_\varepsilon}(m,\omega)$ nor $q_{\beta_\varepsilon}(m)$ depends on $\varepsilon$, so $\Phi_m$ is constant.

For $m \in \{m_1, m_2\}$, $\Phi_m(\varepsilon) = \max_{\hat\sigma} \sum_\omega \bar{u}(\hat\sigma,\omega)\, N_{\beta_\varepsilon}(m,\omega)$ is the pointwise maximum of a family of affine functions of $\varepsilon$, hence \emph{convex} in $\varepsilon$.

The \textbf{direct effect} at each message $m \in \{m_1,m_2\}$ is the contribution holding the optimizer $\hat\sigma$ fixed. The \textbf{indirect effect} is the change in payoff due to the optimizer adjusting.

%% =====================================================================
\section{Envelope Theorem Application}
\label{sec:envelope}
%% =====================================================================

\begin{proposition}[Envelope theorem --- Milgrom--Segal (2002), Theorem~3]
\label{prop:envelope}
Let $f(\varepsilon, x) = \sum_\omega \bar{u}(x, \omega)\, N_{\beta_\varepsilon}(m, \omega)$ where $x$ ranges over a compact set $X$ (the set of private strategies $\hat\sigma$), and define $\Phi_m(\varepsilon) = \max_{x \in X} f(\varepsilon, x)$.

Then $\Phi_m$ has left and right derivatives at $\varepsilon = 0$, given by
\begin{align}
\Phi_m'(0^+) &= \max_{x \in X^*(0)} f_\varepsilon(0, x), \label{eq:env_right}\\
\Phi_m'(0^-) &= \min_{x \in X^*(0)} f_\varepsilon(0, x), \label{eq:env_left}
\end{align}
where $X^*(0) = \argmax_{x \in X} f(0,x)$ is the set of optimizers at $\varepsilon = 0$.
\end{proposition}

\begin{proof}[Verification of hypotheses]
\leavevmode
\begin{enumerate}
\item \textbf{Nonempty argmax:} $X = \Delta(A)^{|\Theta|}$ is compact, $f(0,\cdot)$ is continuous, so $X^*(0) \neq \emptyset$.
\item \textbf{Differentiability in $\varepsilon$:} For each fixed $x$, $\varepsilon \mapsto f(\varepsilon,x)$ is \emph{affine} (because $N_{\beta_\varepsilon}(m,\omega)$ is affine in $\varepsilon$), hence differentiable everywhere.
\item \textbf{Equidifferentiability:} The derivative $f_\varepsilon(\varepsilon, x) = \sum_\omega \bar{u}(x,\omega) \cdot (\pm(1-\alpha)\tau(\mu^*)\mu^*(\omega))$ is \emph{independent of $\varepsilon$} and bounded uniformly in $x$ (because $\bar{u}$ is bounded). Hence the family $\{f(\cdot, x)\}_{x \in X}$ is trivially equidifferentiable.
\end{enumerate}
\end{proof}

\begin{remark}
$\Phi_m(\varepsilon)$ is the pointwise supremum of a family of affine functions of $\varepsilon$, hence convex. (If the action set $A$ is finite, the family is finite and $\Phi_m$ is piecewise linear; for general compact $A$, $\Phi_m$ is still convex and the one-sided derivatives exist by Proposition~\ref{prop:envelope}.) The one-sided derivatives \eqref{eq:env_right}--\eqref{eq:env_left} give the extremal slopes over the argmax set $X^*_m(\beta^*) \coloneqq \argmax_{\hat\sigma \in X} \sum_\omega \bar{u}(\hat\sigma,\omega)\, N_{\beta^*}(m,\omega)$.

When $X^*_m(\beta^*)$ is a singleton (unique Bayes-optimal response), $\Phi_m'(0^+) = \Phi_m'(0^-)$ and $\Phi_m$ is differentiable at $0$. The sketch implicitly assumes this case.
\end{remark}

%% =====================================================================
\section{Selector Family and Corrected FOC}
\label{sec:selector}
%% =====================================================================

This section contains the key repair beyond the original sketch.

\begin{proposition}[Jointly chosen selector family]
\label{prop:selector}
Let $\beta^*$ be a minimizer of $V$ over $\prod_{\mu \in M} \Delta(M)$. Then there exists a Bayes-optimal selector family $(\hat\sigma^*_m)_{m \in M}$ such that for every $\mu \in M$:
\begin{equation}
\label{eq:foc_equal}
\sum_\omega \bar{u}(\hat\sigma^*_{m_1}, \omega)\,\mu(\omega) = \sum_\omega \bar{u}(\hat\sigma^*_{m_2}, \omega)\,\mu(\omega) \quad \forall\, m_1, m_2 \in \supp\,\beta^*(\cdot|\mu),
\end{equation}
and
\begin{equation}
\label{eq:foc_ineq}
\sum_\omega \bar{u}(\hat\sigma^*_m, \omega)\,\mu(\omega) \geq \sum_\omega \bar{u}(\hat\sigma^*_{m'}, \omega)\,\mu(\omega) \quad \forall\, m \in M,\; m' \in \supp\,\beta^*(\cdot|\mu).
\end{equation}
\end{proposition}

\begin{proof}
We proceed in three steps.

\medskip\noindent
\textbf{Step 1: Convexity and subdifferential optimality.}
By \eqref{eq:value} and the definition of $\Phi_m$, $V(\beta) = \sum_m \Phi_m(\beta)$ where each $\Phi_m$ is convex in $\beta$ (being the pointwise maximum of affine functions, as shown in Section~\ref{sec:foc}). Hence $V$ is convex on the compact polyhedron $\prod_\mu \Delta(M)$. Since $\beta^*$ is a minimizer, the subdifferential optimality condition gives
\begin{equation}
\label{eq:subdiff}
0 \in \partial V(\beta^*) + N_{\prod_\mu \Delta(M)}(\beta^*),
\end{equation}
where $N_{\prod_\mu \Delta(M)}(\beta^*)$ is the normal cone to $\prod_\mu \Delta(M)$ at $\beta^*$.

\medskip\noindent
\textbf{Step 2: Subgradient representation.}
Since $V = \sum_m \Phi_m$ and each $\Phi_m(\beta) = \max_{\hat\sigma \in X} \sum_\omega \bar{u}(\hat\sigma,\omega)\, N_\beta(m,\omega)$ is the supremum of affine functions of $\beta$, the subdifferential of $V$ at $\beta^*$ satisfies
\[
\partial V(\beta^*) = \sum_{m \in M} \partial \Phi_m(\beta^*).
\]
For each $m$, the subdifferential of $\Phi_m$ at $\beta^*$ is given by the standard max-of-affine formula (see, e.g., Rockafellar, \textit{Convex Analysis}, Theorem~23.8): it is the convex hull of the gradients of active affine pieces. Concretely, for any $\hat\sigma^*_m \in X^*_m(\beta^*) \coloneqq \argmax_{\hat\sigma} \sum_\omega \bar{u}(\hat\sigma,\omega)\, N_{\beta^*}(m,\omega)$, the vector
\[
g_m(\mu, m') \coloneqq \frac{\partial}{\partial \beta(m'|\mu)} \sum_\omega \bar{u}(\hat\sigma^*_m, \omega)\, N_\beta(m,\omega) \Big|_{\beta = \beta^*}
\]
belongs to $\partial \Phi_m(\beta^*)$. (Here we use the fact that $X$ is convex and $\bar{u}$ is linear in $\hat\sigma$, so any convex combination of active gradients is itself generated by a single mixed strategy in $X^*_m(\beta^*)$.)

Evaluating the partial derivative:
\[
g_m(\mu, m') = \mathbf{1}\{m' = m\} \cdot (1-\alpha)\,\tau(\mu) \sum_\omega \bar{u}(\hat\sigma^*_m, \omega)\,\mu(\omega),
\]
because $N_\beta(m,\omega)$ depends on $\beta(m'|\mu)$ only when $m' = m$ (via the misaligned contribution $(1-\alpha)\tau(\mu)\beta(m|\mu)\mu(\omega)$), and is zero otherwise.

Summing over $m$ produces the subgradient $g \in \partial V(\beta^*)$ with components
\begin{equation}
\label{eq:subgrad_component}
g(\mu, m') = (1-\alpha)\,\tau(\mu) \sum_\omega \bar{u}(\hat\sigma^*_{m'}, \omega)\,\mu(\omega),
\end{equation}
built from the \emph{joint} selector family $(\hat\sigma^*_m)_{m \in M}$. Each $\hat\sigma^*_m$ is Bayes-optimal at $P_{\beta^*}(\cdot|m)$.

\medskip\noindent
\textbf{Step 3: Normal cone condition on $\Delta(M)$.}
The normal cone to $\prod_\mu \Delta(M)$ at $\beta^*$ decomposes row by row. For each $\mu$, the normal cone condition on $\Delta(M)$ at $\beta^*(\cdot|\mu)$ states: there exists $\lambda_\mu \in \R$ such that
\begin{align*}
g_\mu(m) &= \lambda_\mu && \text{if } m \in \supp\,\beta^*(\cdot|\mu), \\
g_\mu(m) &\geq \lambda_\mu && \text{if } m \notin \supp\,\beta^*(\cdot|\mu),
\end{align*}
where $g_\mu(m)$ is the $(\mu,m)$-component of $g$. Unwinding the definition of $g_\mu(m)$ in terms of the selector family, this is exactly the equal-payoff-on-support condition \eqref{eq:foc_equal} (on-support messages achieve the same $\mu$-payoff) and the inequality \eqref{eq:foc_ineq} (off-support messages achieve weakly higher $\mu$-payoff).
\end{proof}

\begin{remark}
The crucial point is that the selector family $(\hat\sigma^*_m)_{m \in M}$ is chosen \emph{jointly} from the subdifferential, not independently at each message. Independent selection would not guarantee the equal-payoff-on-support condition \eqref{eq:foc_equal} for all $\mu$ simultaneously.
\end{remark}

%% =====================================================================
\section{Commitment Game (Block E)}
\label{sec:commitment}
%% =====================================================================

Define the committed agent strategy $\sigma^*$ by $\hat\sigma^*(m) \coloneqq \hat\sigma^*_m$ for each $m \in M$, using the selector family from Proposition~\ref{prop:selector}.

\begin{proposition}[Optimality preservation]
\label{prop:commitment}
$\beta^*$ remains optimal for the adversary in the commitment game where the agent commits to $\sigma^*$.
\end{proposition}

\begin{proof}
Under commitment to $\sigma^*$, the adversary's row problem for each $\mu \in M$ is:
\[
\min_{b \in \Delta(M)} \sum_{m \in M} b(m) \underbrace{\left[\sum_\omega \bar{u}(\hat\sigma^*_m, \omega)\,\mu(\omega)\right]}_{\eqqcolon\, c_m(\mu)}.
\]
This is a linear program over $\Delta(M)$. Its minimizers are precisely the distributions supported on $\argmin_{m \in M} c_m(\mu)$.

By \eqref{eq:foc_equal}, all $m \in \supp\,\beta^*(\cdot|\mu)$ achieve the same value $\kappa_\mu$. By \eqref{eq:foc_ineq}, all $m \notin \supp\,\beta^*(\cdot|\mu)$ achieve $c_m(\mu) \geq \kappa_\mu$. Hence $\supp\,\beta^*(\cdot|\mu) \subseteq \argmin_{m} c_m(\mu)$, so $\beta^*(\cdot|\mu)$ is a minimizer of the row problem. Since the adversary's full objective is the $\tau(\mu)$-weighted sum of row problems, $\beta^*$ minimizes the commitment-game payoff.
\end{proof}

%% =====================================================================
\section{Minimax Conclusion (Block F)}
\label{sec:minimax}
%% =====================================================================

\begin{theorem}[Finite-case alternative proof of Theorem~2]
\label{thm:main}
Suppose $\Omega$, $M = \supp(\tau)$, and $\Theta$ are finite, and $\alpha \in (0,1)$.
\begin{enumerate}[label=(\alph*)]
\item \textbf{(Existence.)} There exist $\sigma^* \in \Sigma$ and $\beta^* \in \mathcal{B}$ such that $(\sigma^*, \beta^*)$ is a saddle point of $U(\beta, \sigma)$:
\[
U(\beta^*, \sigma) \leq U(\beta^*, \sigma^*) \leq U(\beta, \sigma^*) \qquad \forall\, \beta \in \mathcal{B},\; \forall\, \sigma \in \Sigma.
\]
In particular, $\sigma^*$ is robustly rationalizable and optimal, and $\sup_\sigma \inf_\beta U(\beta, \sigma) = \inf_\beta \sup_\sigma U(\beta, \sigma)$.
\item \textbf{(Optimality of any robustly rationalizable strategy.)} If $\sigma$ is any robustly rationalizable strategy, then $\sigma$ is optimal.
\end{enumerate}
\end{theorem}

\begin{proof}
Let $\beta^*$ minimize $V(\beta) = \sup_\sigma U(\beta, \sigma) = \sum_m q_\beta(m)\, W(P_\beta(\cdot|m))$ (the second equality holds by message-by-message Bayes-optimality). Existence of $\beta^*$ follows from compactness of $\prod_\mu \Delta(M)$ and continuity of $V$.

Let $(\hat\sigma^*_m)_{m \in M}$ be the selector family from Proposition~\ref{prop:selector}, and define $\sigma^*$ by $\hat\sigma^*(m) = \hat\sigma^*_m$.

\medskip\noindent
\textbf{Right saddle inequality (F.1):}
By Proposition~\ref{prop:commitment}, $\beta^*$ minimizes the adversary's payoff given $\sigma^*$:
\[
U(\beta^*, \sigma^*) \leq U(\beta, \sigma^*) \qquad \forall\, \beta \in \mathcal{B}.
\]

\textbf{Left saddle inequality (F.2):}
Under $\beta^*$, the posterior at message $m$ is $P_{\beta^*}(\cdot|m)$. By construction, $\hat\sigma^*_m \in \argmax_{\hat\sigma} U(\hat\sigma, P_{\beta^*}(\cdot|m))$. Since $q_{\beta^*}(m) \geq \alpha\,\tau(m) > 0$ for all $m$, every message is on-path, and the payoff is the $q_{\beta^*}(m)$-weighted sum of message-level payoffs. Hence no deviation $\sigma'$ can improve:
\[
U(\beta^*, \sigma') \leq U(\beta^*, \sigma^*) \qquad \forall\, \sigma' \in \Sigma.
\]

Combining (F.1) and (F.2) gives the saddle point. The minimax equality follows by the standard saddle-point theorem. Robust rationalizability of $\sigma^*$ follows because $\beta^*$ is adversarial against $\sigma^*$ and $\hat\sigma^*(m) \in \argmax_{\hat\sigma} U(\hat\sigma, P_{\beta^*}(\cdot|m))$ for all $m$, matching Definition~2 in the paper. This proves part~(a).

\medskip\noindent
\textbf{Part (b): Any robustly rationalizable strategy is optimal.}
Write $U(\sigma) \coloneqq \inf_{\beta \in \mathcal{B}} U(\beta, \sigma)$ for the agent's worst-case payoff. Let $\sigma$ be any robustly rationalizable strategy with adversarial witness $\tilde\beta$. By Definition~2, $\tilde\beta$ is adversarial against $\sigma$ (i.e., $U(\tilde\beta, \sigma) \leq U(\beta, \sigma)$ for all $\beta$), and $\sigma$ is Bayes-optimal message by message under $P_{\tilde\beta}(\cdot|m)$ (i.e., $U(\tilde\beta, \sigma') \leq U(\tilde\beta, \sigma)$ for all $\sigma'$). Together these give
\[
U(\tilde\beta, \sigma') \leq U(\tilde\beta, \sigma) \leq U(\beta, \sigma) \qquad \forall\, \beta,\, \sigma',
\]
so $(\sigma, \tilde\beta)$ is a saddle point and $U(\sigma) = U(\tilde\beta, \sigma) = U(\sigma^*)$ (since all saddle points share the same value). Hence $\sigma$ is optimal.
\end{proof}

%% =====================================================================
\section{Remaining Open Questions}
\label{sec:open}
%% =====================================================================

\begin{enumerate}
\item \textbf{Infinite message space (Block~G).} The FOC approach is inherently finite-dimensional: the perturbation $\beta_\varepsilon(m|\mu^*) = \beta(m|\mu^*) + \varepsilon[\delta_{m_1}(m) - \delta_{m_2}(m)]$ and the subdifferential argument both rely on finite $M$. The author explicitly states: ``I have no idea how to extend the argument.'' This remains open.

\item \textbf{Subdifferential step presentation.} The proof of Proposition~\ref{prop:selector} provides the full convex-analytic argument, but a final manuscript version may benefit from a more streamlined presentation that integrates the subgradient computation with the envelope-theorem output from Proposition~\ref{prop:envelope} more tightly.

\item \textbf{Private types $\theta$.} The proof is written in the simplified setting $u(a,\omega)$ with $|\Theta|=1$. Extension to $|\Theta| > 1$ requires replacing $a^*(\mu)$ with $\hat\sigma^*(\mu)$ and $u(a,\omega)$ with $\bar{u}(\hat\sigma,\omega)$ throughout. The argument goes through unchanged because the envelope theorem applies to the indirect utility $W(\nu) = \max_{\hat\sigma} \sum_\omega \bar{u}(\hat\sigma,\omega)\nu(\omega)$ regardless of how $\hat\sigma$ is parameterized.
\end{enumerate}

\begin{remark}\label{rem:theta}
In the paper's full framework with private types, the Bayes-optimal response is a private strategy $\hat\sigma^*(\mu) \in \Delta(A)^{|\Theta|}$, not a pure action. The envelope theorem still applies because Proposition~\ref{prop:envelope} only requires compactness of the strategy set and affine dependence of the objective on $\varepsilon$, both of which hold for $\hat\sigma$ just as for $a$.
\end{remark}

%% =====================================================================
\section*{Acknowledgments}
%% =====================================================================

This formalization was carried out using the MathPipeProver soft-scaffolding pipeline, with Claude Code (Anthropic) as orchestrator and ChatGPT~5.4 Pro Extended (OpenAI) as the proving engine. Six ChatGPT sessions were used: formalizer (Blocks A--B), prover (Blocks C--D), reviewer (Blocks C--D), prover repair (selector fix), prover (Blocks E--F).

\begin{thebibliography}{99}
\bibitem{DS26} P.~Dworczak and A.~Smolin (2026), ``Robust Trust,'' working paper.
\bibitem{MS02} P.~Milgrom and I.~Segal (2002), ``Envelope Theorems for Arbitrary Choice Sets,'' \textit{Econometrica}, 70(2), 583--601.
\bibitem{Rock70} R.~T.~Rockafellar (1970), \textit{Convex Analysis}, Princeton University Press.
\bibitem{Sion58} M.~Sion (1958), ``On General Minimax Theorems,'' \textit{Pacific Journal of Mathematics}, 8, 171--176.
\end{thebibliography}

\end{document}
